\documentclass[11pt,letterpaper]{article}
\usepackage[utf8]{inputenc}
\usepackage{geometry}
\usepackage{setspace}

\geometry{left=1.0in,right=1.0in,top=0.85in,bottom=0.85in}
\setstretch{1.25}

\setlength{\parindent}{0pt}
\setlength{\parskip}{8pt plus 1pt minus 1pt}

\begin{document}

\thispagestyle{empty}

\begin{center}
    \large\textbf{IN THE UNITED STATES PATENT AND TRADEMARK OFFICE} \\
\end{center}

\vspace{0.2cm}
\noindent\begin{tabular}{p{2.3in}p{4.0in}}
\textbf{In re Patent Application of:} & Reinaldo Javier Aguiar Marcano \\
\textbf{Application Number:} & 19/562,845 \\
\textbf{Confirmation Number:} & 4350 \\
\textbf{Filing Date:} & March 11, 2026 \\
\textbf{Group Art Unit:} & Pre-Exam / TC 2100 \\
\textbf{Examiner:} & Andrea Bennett \\
\textbf{Title of Invention:} & PHYSICALLY DISAGGREGATED COMPUTING SYSTEM \\
\end{tabular}

\vspace{0.6cm}
\begin{center}
    \textbf{\large REPLACEMENT ABSTRACT (37 CFR \S{} 1.72(b))} \\
    \textbf{(COMMENCING ON A SEPARATE SHEET)}
\end{center}

\vspace{0.4cm}

\begin{center}
    \textbf{ABSTRACT}
\end{center}

A computing system architecture where primary computer components are physically disaggregated into separate enclosures and interconnected by an ultra-high-bandwidth, low-latency data fabric. The system comprises a first chassis containing a plurality of processors, a second chassis containing a plurality of memory modules, and potentially other chassis for storage and accelerator cards. The data fabric provides communication between the disaggregated components at speeds sufficient to allow them to operate as a single, cohesive logical computer, effectively 'stretching' the motherboard across a data center. This architecture overcomes the thermal, power, and resource-stranding inefficiencies of traditional monolithic server designs.

\end{document}
